\subsection{Transition Rules}
{\small\renewcommand{\arraystretch}{1.09}
\begin{tabularx}{\linewidth}{@{}p{2.65cm}p{4.25cm}Y@{}}
\toprule
Current state & Event or condition & Result\\
\midrule
Standby & Valid coin; credit below 50 & Add credit; stay in Standby.\\
Standby & Valid coin; credit reaches 50 & Add credit; enter Ready.\\
Ready & Valid coin & Add credit; stay in Ready.\\
Ready & RUN & Clear credit, save start time and enter Running.\\
Running & PAUSE & Turn washing off; enter Paused. Keep start time.\\
Paused & RUN & Turn washing on; enter Running. Keep start time.\\
Running or Paused & First STOP & Set STOP count to 1. Keep the current state.\\
Running or Paused & Second STOP & Clear cycle data; enter Standby.\\
Running or Paused & Elapsed time reaches 30 min & Clear cycle data; enter Standby.\\
Any normal state & Fault signal & Turn washing off, clear credit and enter Error.\\
Error & Board reset after fault is cleared & Return to Standby with zero credit.\\
\bottomrule
\end{tabularx}}
Unsupported coins do not add credit. Other unused controls leave the state unchanged.
RUN in Running does not restart the timer.

\subsection{Timer and Pause Behavior}
The timer uses elapsed time from the first RUN:
\[
T_{\mathrm{cycle}}=30\times60\times1000=1{,}800{,}000\ \mathrm{ms},\qquad
T_{\mathrm{left}}=\max(0,T_{\mathrm{cycle}}-(t-t_{\mathrm{start}})).
\]
The start time changes only when a paid cycle starts. For example, pause at
minute 8 and resume at minute 16. There are 14 minutes left after resuming, and
the cycle still ends at minute 30.
\fig{04_pause_timeline.pdf}{1.0}{Example cycle with eight minutes spent in Paused.}

\subsection{Work in Each Loop}
\begin{enumerate}
\item Read the inputs and create debounced press events.
\item Check the fault signal. A fault has the highest priority.
\item Check the deadline in both Running and Paused.
\item Collect valid coins before a cycle, or process STOP, PAUSE and RUN during a cycle.
\item Update the LEDs and wash output from the current state.
\item Write a short trace when the state, credit or STOP count changes.
\end{enumerate}
In an active cycle, STOP is processed before the other buttons. PAUSE takes
priority over RUN when the machine is Running. An expired cycle ignores input
events from that same loop. The loop uses no blocking wait.
